\begin{proofbox}
	\textbf{\textcolor{CProof}{思路提示}}

	由已知等式$f(x+a)=f(x-a)$，通过变量代换$x\mapsto x+a$，可直接构造出$f(x+2a)=f(x)$，从而判定周期$2a$。

	\medskip
	\textbf{\textcolor{CProof}{正式推导}}
	\begin{description}[style=nextline, leftmargin=2.8em, labelsep=0.5em, itemsep=0.55em]
		\item[\textbf{步骤一（条件陈列）：}] 已知对任意$x$，有
			\[
				f(x+a)=f(x-a).
			\]
			\par\textbf{理由：} 这是题设直接给出的条件。

		\item[\textbf{步骤二（变量替换）：}] 将上式中的$x$替换为$x+a$（因为条件对任意$x$成立，所以对$x+a$也成立）：
			\[
				f((x+a)+a)=f((x+a)-a).
			\]
			\par\textbf{理由：} 条件中$x$的任意性允许我们代入任何实数。

		\item[\textbf{步骤三（化简等式）：}] 化简括号：
			\[
				f(x+2a)=f(x).
			\]
			\par\textbf{理由：} $(x+a)+a=x+2a$，$(x+a)-a=x$。

		\item[\textbf{步骤四（周期判定）：}] 对照周期定义$f(x+T)=f(x)$，取$T=2a$即满足。因此$2a$是$f(x)$的一个周期。
			\par\textbf{理由：} 周期函数要求存在非零常数$T$使得$f(x+T)=f(x)$对任意$x$成立，这里$T=2a$满足且$a\neq 0$保证$T\neq 0$。
	\end{description}

	\medskip
	\textbf{\textcolor{CProof}{结论回扣}}

	\textbf{综上，由$f(x+a)=f(x-a)$（$a\neq 0$）可推出$f(x+2a)=f(x)$，即周期$T=2a$。这一推导的关键步骤是用$x+a$代替$x$。}
\end{proofbox}
